\documentclass[11pt]{article}

\usepackage[margin=1in]{geometry}
\usepackage{amsmath,amssymb}
\usepackage{booktabs}
\usepackage{array}
\usepackage{hyperref}

\title{Understanding the MPC Projection Error}
\author{}
\date{September 1, 2026}

\begin{document}
\maketitle

\section{The central issue}

Using Cartesian coordinates is not itself an error, and using Frenet
coordinates is not the only valid design. The earlier local MPC became
inconsistent because it predicted the vehicle in Cartesian coordinates but
evaluated that prediction through an approximate Frenet coordinate system tied
to a preselected future reference point. The approximation was then also used
in tracking and safety constraints as though it were the predicted point's true
Frenet projection.

There are four formulations to distinguish:

\begin{enumerate}
  \item the Frenet positional formulation presented in the
  $\alpha$-RACER paper;
  \item the world-coordinate MPC used by the released project code;
  \item the earlier local hybrid implementation containing the projection
  error; and
  \item the current local four-state Frenet MPC.
\end{enumerate}

\section{What a true Frenet projection means}

Let the track centreline be the curve
\[
  C(s)=
  \begin{bmatrix}
    x_c(s)\\y_c(s)
  \end{bmatrix},
\]
with unit tangent $t(s)$ and unit normal $n(s)$. For a vehicle position in
world coordinates,
\[
  X=
  \begin{bmatrix}
    x\\y
  \end{bmatrix},
\]
the longitudinal Frenet coordinate is obtained by finding the closest point on
the centreline:
\[
  s^*=\underset{s}{\operatorname{arg\,min}}
  \left\|X-C(s)\right\|_2^2.
\]
The signed lateral coordinate is then
\[
  e^*=n(s^*)^\top\bigl(X-C(s^*)\bigr).
\]
Thus the projection
\[
  X\longmapsto(s^*,e^*)
\]
depends on the position being projected. The correct centreline section cannot
in general be selected before the future position $X$ is known.

\section{The earlier local hybrid MPC}

The earlier local MPC predicted a Cartesian kinematic state
\[
  X_k=(x_k,y_k,\psi_k,v_k)
\]
while the high-level policy supplied a Frenet reference
\[
  p_k^{\mathrm{ref}}
  =
  (s_k^{\mathrm{ref}},e_k^{\mathrm{ref}}).
\]
To compare these two representations, the implementation preselected the
centreline section associated with the $k$th reference point:
\[
  C_k=C(s_k^{\mathrm{ref}}),\qquad
  t_k=t(s_k^{\mathrm{ref}}),\qquad
  n_k=n(s_k^{\mathrm{ref}}).
\]
It then assigned the symbolic Cartesian prediction the approximate coordinates
\[
  \widetilde s_k
  =s_k^{\mathrm{ref}}
  +t_k^\top(X_k^{xy}-C_k),
\]
\[
  \widetilde e_k
  =n_k^\top(X_k^{xy}-C_k),
\]
where $X_k^{xy}=(x_k,y_k)^\top$. Tracking and boundary costs were evaluated
using $(\widetilde s_k,\widetilde e_k)$.

This is a local tangent-line approximation around the reference section. It is
reasonable only when
\[
  X_k^{xy}\approx C(s_k^{\mathrm{ref}}).
\]
The actual projection instead requires
\[
  s_k^*=\underset{s}{\operatorname{arg\,min}}
  \|X_k^{xy}-C(s)\|_2^2,
\]
followed by $n(s_k^*)$, which need not equal the preselected
$n(s_k^{\mathrm{ref}})$.

\subsection{Why the error becomes important}

The approximation can deteriorate when the predicted vehicle:

\begin{itemize}
  \item travels more slowly than its time-advanced reference;
  \item is delayed by another vehicle;
  \item moves away from the reference near a sharp bend;
  \item approaches a track or collision constraint; or
  \item becomes nearly stationary while the reference horizon continues to
  advance.
\end{itemize}

In a hairpin, two nearby centreline sections may have substantially different
normal directions. A primarily longitudinal lag can therefore appear as a
large lateral error when measured in the wrong section. If a vehicle stops
while $s_k^{\mathrm{ref}}$ continues to advance, later predictions are measured
in increasingly distant local frames. This can exaggerate predicted lateral
cost and make acceleration appear less attractive.

The projection error could therefore reinforce the observed low-speed local
solution. It was not, however, the only cause of that behavior.

\subsection{Why exact projection inside the old NLP is difficult}

One could in principle define the true projection at every predicted step as
\[
  s_k^*=\underset{s}{\operatorname{arg\,min}}
  \|X_k^{xy}-C(s)\|_2^2.
\]
However, $X_k^{xy}$ is itself an optimization variable. A polyline-based exact
projection would introduce nearest-segment selection, discontinuous changes of
the selected segment, start/finish wrap-around, and an inner minimization into
the nonlinear program. A smooth implicit projection is possible, but it would
make the solver considerably more complicated. The cleaner solution is to use
a coordinate formulation that does not require this repeated symbolic
projection.

\section{What the released project actually does}

The released project avoids this hybrid projection. Its MPC state is described
in some comments as $[s,e,\psi,v]$, but the implemented dynamics and runtime
inputs are Cartesian:
\[
\begin{aligned}
  x_{k+1} &= x_k+v_k\cos\psi_k\,\Delta t,\\
  y_{k+1} &= y_k+v_k\sin\psi_k\,\Delta t,\\
  \psi_{k+1}
  &=\psi_k+\frac{v_k}{L}\tan(S_a u_k^\delta)\,\Delta t,\\
  v_{k+1} &=v_k+a_k\Delta t.
\end{aligned}
\]
The controller receives the initial state $[x,y,\psi,v]$ and minimizes the
world-coordinate position error
\[
  J_{\mathrm{track}}
  =\sum_k
  \left\|
    \begin{bmatrix}x_k\\y_k\end{bmatrix}
    -
    \begin{bmatrix}x_k^{\mathrm{ref}}\\y_k^{\mathrm{ref}}\end{bmatrix}
  \right\|_Q^2.
\]

Before calling the MPC, the race code generates a reference trajectory outside
the optimizer. It advances along the raceline using the speed profile, computes
the overtaking and blocking offsets, converts the shifted points to world
coordinates, and sends those world points to the MPC. Its data flow is
\[
  \text{Frenet-style tactical offsets}
  \longrightarrow
  \text{world-coordinate reference}
  \longrightarrow
  \text{world-coordinate MPC}.
\]
This formulation is internally consistent and requires no symbolic
world-to-Frenet projection.

It is also less constrained than the current local MPC. The released MPC does
not contain the same explicit Frenet track-boundary and smooth collision
constraints. The numerical simulator instead applies a separate recovery after
$|e|>0.55$ m. Thus its simpler coordinate formulation should not be interpreted
as providing the same constrained optimization problem.

The relevant released files are:
\begin{itemize}
  \item \url{https://github.com/sastry-group/alpha-RACER/blob/main/car_ros2/car_ros2/mpc_controller.py};
  \item \url{https://github.com/sastry-group/alpha-RACER/blob/main/car_ros2/car_ros2/multi_car_comp_blocking.py}.
\end{itemize}

\section{The paper-level Frenet formulation}

The $\alpha$-RACER paper separates the vehicle dynamics state from its
track-relative position. The vehicle state $x$ may retain the quantities needed
by the dynamic bicycle model, while
\[
  p_k=(p_{x,k},p_{y,k})
\]
represents longitudinal progress and lateral position in the Frenet frame. The
optimization propagates both the vehicle dynamics and the track-relative
position:
\[
  x_{k+1}=f(x_k,u_k),
  \qquad
  p_{k+1}=g(p_k,x_k,u_k,\kappa).
\]
Its positional objective can then be written directly as
\[
  J_{\mathrm{track}}
  =\sum_k
  \left\|p_k-p_k^{\mathrm{ref}}\right\|_Q^2.
\]
Track containment and interaction constraints are also expressed using the
Frenet components, for example
\[
  |p_{y,k}|\leq\frac{w_{\max}}{2}.
\]

The important idea is not that every physical state must be discarded in favor
of two Frenet variables. Rather, the Frenet position $p_k$ is part of the
optimization itself. It is not reconstructed from each symbolic Cartesian
prediction using a fixed future reference section. The paper is available at
\url{https://arxiv.org/pdf/2412.08855v3}.

The released project code and the paper formulation are therefore not identical:
the paper presents a Frenet positional optimization, whereas the released MPC
implements the simpler world-reference tracker described above.

\section{The current local Frenet MPC}

The current local MPC uses the reduced internal state
\[
  z_k=
  \begin{bmatrix}
    s_k&e_k&\Delta\psi_k&v_k
  \end{bmatrix}^{\top}.
\]
The measured Cartesian state is projected once at the beginning of each solve:
\[
  (x_0,y_0,\psi_0,v_0)
  \longrightarrow
  (s_0,e_0,\Delta\psi_0,v_0).
\]
The optimizer then propagates Frenet dynamics directly:
\[
  s_{k+1}
  =s_k+\Delta t
  \frac{v_k\cos\Delta\psi_k}{1-\kappa(s_k)e_k},
\]
\[
  e_{k+1}
  =e_k+\Delta t\,v_k\sin\Delta\psi_k,
\]
\[
  \Delta\psi_{k+1}
  =\Delta\psi_k+\Delta t
  \left[
    \frac{v_k}{L}\tan\delta_k
    -\kappa(s_k)
     \frac{v_k\cos\Delta\psi_k}{1-\kappa(s_k)e_k}
  \right].
\]
The track curvature $\kappa(s_k)$ is evaluated inside the NLP through a smooth
periodic interpolation. The tracking cost is now directly
\[
  J_{\mathrm{track}}
  =\sum_k
  \left\|
    \begin{bmatrix}s_k\\e_k\end{bmatrix}
    -
    \begin{bmatrix}s_k^{\mathrm{ref}}\\e_k^{\mathrm{ref}}\end{bmatrix}
  \right\|_Q^2.
\]
Likewise, the boundary is evaluated from $e_k$ itself, and collision separation
is evaluated from the predicted Frenet differences
\[
  \Delta s_k=s_k-s_k^{\mathrm{opp}},
  \qquad
  \Delta e_k=e_k-e_k^{\mathrm{opp}}.
\]
No future Cartesian prediction needs to be projected back into Frenet
coordinates.

\section{Comparison}

\begin{center}
\small
\renewcommand{\arraystretch}{1.25}
\begin{tabular}{>{\raggedright\arraybackslash}p{0.19\textwidth}
                >{\raggedright\arraybackslash}p{0.22\textwidth}
                >{\raggedright\arraybackslash}p{0.23\textwidth}
                >{\raggedright\arraybackslash}p{0.25\textwidth}}
\toprule
Formulation & Prediction variables & Reference representation & Tracking error \\
\midrule
Paper & Physical state $x$ plus Frenet position $p$ & Frenet $p^{\mathrm{ref}}$ & Directly compares $p_k-p_k^{\mathrm{ref}}$ \\
Released MPC & Cartesian $(x,y,\psi,v)$ & Precomputed world trajectory & Direct world-to-world comparison \\
Earlier local MPC & Cartesian $(x,y,\psi,v)$ & Frenet reference & Fixed future cross-section used as an approximate projection \\
Current local MPC & Frenet $(s,e,\Delta\psi,v)$ & Frenet reference & Direct Frenet-to-Frenet comparison \\
\bottomrule
\end{tabular}
\end{center}

\section{What the correction established}

After the Frenet rewrite, a 25 s three-car $\alpha$-RACER screen reached
\[
  106.12/109.69/106.27\ \mathrm{m}
\]
with no boundary violation, recovery activation, collision, or failed MPC
solve. This supports the conclusion that the hybrid projection was a real local
formulation error.

Parameter-IBR could nevertheless still stall near the same hairpin. The
remaining failure therefore cannot be attributed only to projection. It also
depends on the high-level reference selected by Parameter-IBR, the finite MPC
horizon, soft safety costs, and the low-speed controllability of the kinematic
model.

The concise conclusion is:

\begin{quote}
The released MPC uses a consistent but relatively permissive world-to-world
tracking formulation. The paper presents a Frenet positional formulation. The
earlier local implementation mixed Cartesian predictions with a fixed-section
Frenet approximation, which was incorrect when the prediction departed from
the selected reference section. The current local MPC removes this
approximation by predicting directly in four-state Frenet coordinates.
\end{quote}

\end{document}
